% Options for packages loaded elsewhere
\PassOptionsToPackage{unicode}{hyperref}
\PassOptionsToPackage{hyphens}{url}
\PassOptionsToPackage{dvipsnames,svgnames,x11names}{xcolor}
%
\documentclass[
  letterpaper,
  DIV=11,
  numbers=noendperiod]{scrartcl}

\usepackage{amsmath,amssymb}
\usepackage{iftex}
\ifPDFTeX
  \usepackage[T1]{fontenc}
  \usepackage[utf8]{inputenc}
  \usepackage{textcomp} % provide euro and other symbols
\else % if luatex or xetex
  \usepackage{unicode-math}
  \defaultfontfeatures{Scale=MatchLowercase}
  \defaultfontfeatures[\rmfamily]{Ligatures=TeX,Scale=1}
\fi
\usepackage{lmodern}
\ifPDFTeX\else  
    % xetex/luatex font selection
\fi
% Use upquote if available, for straight quotes in verbatim environments
\IfFileExists{upquote.sty}{\usepackage{upquote}}{}
\IfFileExists{microtype.sty}{% use microtype if available
  \usepackage[]{microtype}
  \UseMicrotypeSet[protrusion]{basicmath} % disable protrusion for tt fonts
}{}
\makeatletter
\@ifundefined{KOMAClassName}{% if non-KOMA class
  \IfFileExists{parskip.sty}{%
    \usepackage{parskip}
  }{% else
    \setlength{\parindent}{0pt}
    \setlength{\parskip}{6pt plus 2pt minus 1pt}}
}{% if KOMA class
  \KOMAoptions{parskip=half}}
\makeatother
\usepackage{xcolor}
\setlength{\emergencystretch}{3em} % prevent overfull lines
\setcounter{secnumdepth}{-\maxdimen} % remove section numbering
% Make \paragraph and \subparagraph free-standing
\makeatletter
\ifx\paragraph\undefined\else
  \let\oldparagraph\paragraph
  \renewcommand{\paragraph}{
    \@ifstar
      \xxxParagraphStar
      \xxxParagraphNoStar
  }
  \newcommand{\xxxParagraphStar}[1]{\oldparagraph*{#1}\mbox{}}
  \newcommand{\xxxParagraphNoStar}[1]{\oldparagraph{#1}\mbox{}}
\fi
\ifx\subparagraph\undefined\else
  \let\oldsubparagraph\subparagraph
  \renewcommand{\subparagraph}{
    \@ifstar
      \xxxSubParagraphStar
      \xxxSubParagraphNoStar
  }
  \newcommand{\xxxSubParagraphStar}[1]{\oldsubparagraph*{#1}\mbox{}}
  \newcommand{\xxxSubParagraphNoStar}[1]{\oldsubparagraph{#1}\mbox{}}
\fi
\makeatother

\usepackage{color}
\usepackage{fancyvrb}
\newcommand{\VerbBar}{|}
\newcommand{\VERB}{\Verb[commandchars=\\\{\}]}
\DefineVerbatimEnvironment{Highlighting}{Verbatim}{commandchars=\\\{\}}
% Add ',fontsize=\small' for more characters per line
\usepackage{framed}
\definecolor{shadecolor}{RGB}{241,243,245}
\newenvironment{Shaded}{\begin{snugshade}}{\end{snugshade}}
\newcommand{\AlertTok}[1]{\textcolor[rgb]{0.68,0.00,0.00}{#1}}
\newcommand{\AnnotationTok}[1]{\textcolor[rgb]{0.37,0.37,0.37}{#1}}
\newcommand{\AttributeTok}[1]{\textcolor[rgb]{0.40,0.45,0.13}{#1}}
\newcommand{\BaseNTok}[1]{\textcolor[rgb]{0.68,0.00,0.00}{#1}}
\newcommand{\BuiltInTok}[1]{\textcolor[rgb]{0.00,0.23,0.31}{#1}}
\newcommand{\CharTok}[1]{\textcolor[rgb]{0.13,0.47,0.30}{#1}}
\newcommand{\CommentTok}[1]{\textcolor[rgb]{0.37,0.37,0.37}{#1}}
\newcommand{\CommentVarTok}[1]{\textcolor[rgb]{0.37,0.37,0.37}{\textit{#1}}}
\newcommand{\ConstantTok}[1]{\textcolor[rgb]{0.56,0.35,0.01}{#1}}
\newcommand{\ControlFlowTok}[1]{\textcolor[rgb]{0.00,0.23,0.31}{\textbf{#1}}}
\newcommand{\DataTypeTok}[1]{\textcolor[rgb]{0.68,0.00,0.00}{#1}}
\newcommand{\DecValTok}[1]{\textcolor[rgb]{0.68,0.00,0.00}{#1}}
\newcommand{\DocumentationTok}[1]{\textcolor[rgb]{0.37,0.37,0.37}{\textit{#1}}}
\newcommand{\ErrorTok}[1]{\textcolor[rgb]{0.68,0.00,0.00}{#1}}
\newcommand{\ExtensionTok}[1]{\textcolor[rgb]{0.00,0.23,0.31}{#1}}
\newcommand{\FloatTok}[1]{\textcolor[rgb]{0.68,0.00,0.00}{#1}}
\newcommand{\FunctionTok}[1]{\textcolor[rgb]{0.28,0.35,0.67}{#1}}
\newcommand{\ImportTok}[1]{\textcolor[rgb]{0.00,0.46,0.62}{#1}}
\newcommand{\InformationTok}[1]{\textcolor[rgb]{0.37,0.37,0.37}{#1}}
\newcommand{\KeywordTok}[1]{\textcolor[rgb]{0.00,0.23,0.31}{\textbf{#1}}}
\newcommand{\NormalTok}[1]{\textcolor[rgb]{0.00,0.23,0.31}{#1}}
\newcommand{\OperatorTok}[1]{\textcolor[rgb]{0.37,0.37,0.37}{#1}}
\newcommand{\OtherTok}[1]{\textcolor[rgb]{0.00,0.23,0.31}{#1}}
\newcommand{\PreprocessorTok}[1]{\textcolor[rgb]{0.68,0.00,0.00}{#1}}
\newcommand{\RegionMarkerTok}[1]{\textcolor[rgb]{0.00,0.23,0.31}{#1}}
\newcommand{\SpecialCharTok}[1]{\textcolor[rgb]{0.37,0.37,0.37}{#1}}
\newcommand{\SpecialStringTok}[1]{\textcolor[rgb]{0.13,0.47,0.30}{#1}}
\newcommand{\StringTok}[1]{\textcolor[rgb]{0.13,0.47,0.30}{#1}}
\newcommand{\VariableTok}[1]{\textcolor[rgb]{0.07,0.07,0.07}{#1}}
\newcommand{\VerbatimStringTok}[1]{\textcolor[rgb]{0.13,0.47,0.30}{#1}}
\newcommand{\WarningTok}[1]{\textcolor[rgb]{0.37,0.37,0.37}{\textit{#1}}}

\providecommand{\tightlist}{%
  \setlength{\itemsep}{0pt}\setlength{\parskip}{0pt}}\usepackage{longtable,booktabs,array}
\usepackage{calc} % for calculating minipage widths
% Correct order of tables after \paragraph or \subparagraph
\usepackage{etoolbox}
\makeatletter
\patchcmd\longtable{\par}{\if@noskipsec\mbox{}\fi\par}{}{}
\makeatother
% Allow footnotes in longtable head/foot
\IfFileExists{footnotehyper.sty}{\usepackage{footnotehyper}}{\usepackage{footnote}}
\makesavenoteenv{longtable}
\usepackage{graphicx}
\makeatletter
\newsavebox\pandoc@box
\newcommand*\pandocbounded[1]{% scales image to fit in text height/width
  \sbox\pandoc@box{#1}%
  \Gscale@div\@tempa{\textheight}{\dimexpr\ht\pandoc@box+\dp\pandoc@box\relax}%
  \Gscale@div\@tempb{\linewidth}{\wd\pandoc@box}%
  \ifdim\@tempb\p@<\@tempa\p@\let\@tempa\@tempb\fi% select the smaller of both
  \ifdim\@tempa\p@<\p@\scalebox{\@tempa}{\usebox\pandoc@box}%
  \else\usebox{\pandoc@box}%
  \fi%
}
% Set default figure placement to htbp
\def\fps@figure{htbp}
\makeatother

\KOMAoption{captions}{tableheading}
\makeatletter
\@ifpackageloaded{caption}{}{\usepackage{caption}}
\AtBeginDocument{%
\ifdefined\contentsname
  \renewcommand*\contentsname{Table of contents}
\else
  \newcommand\contentsname{Table of contents}
\fi
\ifdefined\listfigurename
  \renewcommand*\listfigurename{List of Figures}
\else
  \newcommand\listfigurename{List of Figures}
\fi
\ifdefined\listtablename
  \renewcommand*\listtablename{List of Tables}
\else
  \newcommand\listtablename{List of Tables}
\fi
\ifdefined\figurename
  \renewcommand*\figurename{Figure}
\else
  \newcommand\figurename{Figure}
\fi
\ifdefined\tablename
  \renewcommand*\tablename{Table}
\else
  \newcommand\tablename{Table}
\fi
}
\@ifpackageloaded{float}{}{\usepackage{float}}
\floatstyle{ruled}
\@ifundefined{c@chapter}{\newfloat{codelisting}{h}{lop}}{\newfloat{codelisting}{h}{lop}[chapter]}
\floatname{codelisting}{Listing}
\newcommand*\listoflistings{\listof{codelisting}{List of Listings}}
\makeatother
\makeatletter
\makeatother
\makeatletter
\@ifpackageloaded{caption}{}{\usepackage{caption}}
\@ifpackageloaded{subcaption}{}{\usepackage{subcaption}}
\makeatother

\usepackage{bookmark}

\IfFileExists{xurl.sty}{\usepackage{xurl}}{} % add URL line breaks if available
\urlstyle{same} % disable monospaced font for URLs
\hypersetup{
  pdftitle={Process 30 Day Map Challenge},
  pdfauthor={Malcolm Campbell},
  colorlinks=true,
  linkcolor={blue},
  filecolor={Maroon},
  citecolor={Blue},
  urlcolor={Blue},
  pdfcreator={LaTeX via pandoc}}


\title{Process 30 Day Map Challenge}
\author{Malcolm Campbell}
\date{}

\begin{document}
\maketitle


\section{Process: 30 Day Map
Challenge}\label{process-30-day-map-challenge}

This document is for Day 23. Process. Show how you make a map. This
could be a tutorial, a step-by-step graphic, a blog post, a video, or a
screenshot of your work environment. Combine it with a map from another
day! See the full list here {[}30 Day Map Challenge{]}
\url{https://30daymapchallenge.com/\#data}

\subsection{Making a `simple' map}\label{making-a-simple-map}

To make a figure of the North South Divide in the UK, we will need to
find some administrative boundaries from an authoritative source.
Firstly, lets load some key packages we will need using the
\texttt{pacman} library and the \texttt{p\_load} function which can help
us by installing any library that is missing. We will want to use
rnaturalearth - for the boundaries with sf and tmap - for the spatial
elements and visualisation alongside tidyverse - for data wrangling.

\begin{Shaded}
\begin{Highlighting}[]
\NormalTok{pacman}\SpecialCharTok{::}\FunctionTok{p\_load}\NormalTok{(rnaturalearth, sf, tmap, tidyverse)}
\end{Highlighting}
\end{Shaded}

Next, lets get the administrative boundaries of the United Kingdom to
use in our figure. To do this we use the ne\_states() command from
\{rnaturalearth\} library which gives us sub national boundaries. We
need to input the country required, for example the United Kingdom, the
file format we want back, in this case an \texttt{sf} format.

\begin{Shaded}
\begin{Highlighting}[]
\NormalTok{UK }\OtherTok{\textless{}{-}} \FunctionTok{ne\_states}\NormalTok{(}\AttributeTok{country=}\StringTok{"United Kingdom"}\NormalTok{, }\AttributeTok{returnclass =} \StringTok{"sf"}\NormalTok{)}
\end{Highlighting}
\end{Shaded}

A quick test to see if this has downloaded as desired is to produce a
simple map using \texttt{tmap}, or whichever spatial visualisation works
for you. Other possibilities include \texttt{ggplot2} and
\texttt{leaflet} among others. For \texttt{tmap}, we need to specify the
file (UK) within the \texttt{tm\_shape} function and then use either
\texttt{tm\_polygons}, \texttt{tm\_borders} or \texttt{tm\_fill} to
create a map. Have a look at the three maps below to see the subtle
differences and then query \texttt{?tm\_borders} using the question mark
to understand the arguments needed for each.

\begin{Shaded}
\begin{Highlighting}[]
\FunctionTok{tm\_shape}\NormalTok{(UK) }\SpecialCharTok{+} \FunctionTok{tm\_polygons}\NormalTok{()}
\end{Highlighting}
\end{Shaded}

\pandocbounded{\includegraphics[keepaspectratio]{Day23_Lesson_files/figure-pdf/map UK-1.pdf}}

\begin{Shaded}
\begin{Highlighting}[]
\FunctionTok{tm\_shape}\NormalTok{(UK) }\SpecialCharTok{+} \FunctionTok{tm\_borders}\NormalTok{()}
\end{Highlighting}
\end{Shaded}

\pandocbounded{\includegraphics[keepaspectratio]{Day23_Lesson_files/figure-pdf/map UK-2.pdf}}

\begin{Shaded}
\begin{Highlighting}[]
\FunctionTok{tm\_shape}\NormalTok{(UK) }\SpecialCharTok{+} \FunctionTok{tm\_fill}\NormalTok{()}
\end{Highlighting}
\end{Shaded}

\pandocbounded{\includegraphics[keepaspectratio]{Day23_Lesson_files/figure-pdf/map UK-3.pdf}}

Next, lets decide which places are in the North, and which are in the
south. We can do this by seeing what regions we have available to put
into categories.

\begin{Shaded}
\begin{Highlighting}[]
\FunctionTok{unique}\NormalTok{(UK}\SpecialCharTok{$}\NormalTok{region)}
\end{Highlighting}
\end{Shaded}

\begin{verbatim}
 [1] "Northern Ireland"           "West Wales and the Valleys"
 [3] "North West"                 "West Midlands"             
 [5] "South West"                 "Eastern"                   
 [7] "North East"                 "South Western"             
 [9] "Yorkshire and the Humber"   "East Midlands"             
[11] "East"                       "South East"                
[13] "East Wales"                 "Highlands and Islands"     
[15] "North Eastern"              "Greater London"            
\end{verbatim}

After seeing that there are 16 regions, we want to assign a value to
each in order to be able to map the North-South divide. Using the
\texttt{tidyverse} package to do some basic data wrangling and variable
creation with the \texttt{mutate} and the \texttt{case\_when} functions
will help us get there. We will create a variable called \texttt{DIVIDE}
to use as a category for future mapping by creating a binary north or
south set of responses for each region in the spatial file of the UK.

\begin{Shaded}
\begin{Highlighting}[]
\NormalTok{UK }\OtherTok{\textless{}{-}}\NormalTok{ UK }\SpecialCharTok{|\textgreater{}}
  \FunctionTok{mutate}\NormalTok{ ( }\AttributeTok{DIVIDE =} \FunctionTok{case\_when}\NormalTok{(}
\NormalTok{    region }\SpecialCharTok{==}\StringTok{"South East"} \SpecialCharTok{\textasciitilde{}} \StringTok{"South"}\NormalTok{,}
\NormalTok{    region }\SpecialCharTok{==}\StringTok{"East"} \SpecialCharTok{\textasciitilde{}} \StringTok{"South"}\NormalTok{,}
\NormalTok{    region }\SpecialCharTok{==}\StringTok{"South West"} \SpecialCharTok{\textasciitilde{}} \StringTok{"South"}\NormalTok{,}
\NormalTok{    region }\SpecialCharTok{==}\StringTok{"Greater London"} \SpecialCharTok{\textasciitilde{}} \StringTok{"South"}\NormalTok{,}
    \ConstantTok{TRUE} \SpecialCharTok{\textasciitilde{}} \StringTok{"North"}\NormalTok{ )) }\CommentTok{\# Default category {-} for the other 12}
\end{Highlighting}
\end{Shaded}

Now for the slightly more complicated part. We can now set up a custom
palette to use in mapping. To do this we select our desired colours for
the north and the south.

\begin{Shaded}
\begin{Highlighting}[]
\NormalTok{NS\_palette }\OtherTok{\textless{}{-}} \FunctionTok{c}\NormalTok{(}\StringTok{"South"} \OtherTok{=}\StringTok{"dodgerblue2"}\NormalTok{, }\AttributeTok{North=}\StringTok{"firebrick"}\NormalTok{)}
\end{Highlighting}
\end{Shaded}

Next, we will make a map using the \texttt{tmap} package. It is easier
to do this in steps, to try and ease understanding of the different
commands and options, which are myriad!

Firstly, lets just colour in map polygons and add a title using
\texttt{tm\_polygons} and \texttt{tm\_title} commands. A title is fairly
easy to explain, we simply add the text ``North-South Divide'' within
the command so that
\texttt{tm\_title(text="North-South\ Divide",\ size\ =\ 2.5)} becomes
the title and we can adjust the size also, lets go for 2.5 times bigger.

We use our shape, the UK, within the \texttt{tm\_shape} command as
before. This time however, we want to fill in the polygons using the
DIVIDE variable, \texttt{fill="DIVIDE"} that we created. The critical
part is to use \texttt{fill.scale} to set the values from our previously
created custom palette. To do this we use
\texttt{tm\_scale(values\ =\ NS\_palette)} which will pull either the
north or south colours that we set up in our custom palette and assign
these to each region in the map. Finally, we set the border of the
polygon colours to NULL \texttt{col=NULL}, as we haven't properly merged
the underlying polygons and we want to hide this. The last step is then
to suppress the legend using the \texttt{tm\_legend\_hide()} option.
I've also added a command to stop tmap autoscaling using
\texttt{tmap\_options(component.autoscale\ =\ FALSE)}.

\begin{Shaded}
\begin{Highlighting}[]
\NormalTok{NSmap }\OtherTok{\textless{}{-}}
  \FunctionTok{tm\_shape}\NormalTok{(UK) }\SpecialCharTok{+}
  \FunctionTok{tm\_polygons}\NormalTok{(}\AttributeTok{fill=}\StringTok{"DIVIDE"}\NormalTok{,}
              \AttributeTok{fill.scale =} \FunctionTok{tm\_scale}\NormalTok{(}\AttributeTok{values =}\NormalTok{ NS\_palette),}
              \AttributeTok{col=}\ConstantTok{NULL}\NormalTok{, }\AttributeTok{fill.legend =} \FunctionTok{tm\_legend\_hide}\NormalTok{())}\SpecialCharTok{+}
  \FunctionTok{tm\_title}\NormalTok{(}\AttributeTok{text=}\StringTok{"North{-}South Divide"}\NormalTok{, }\AttributeTok{size =} \FloatTok{2.5}\NormalTok{) }\SpecialCharTok{+}
  \FunctionTok{tmap\_options}\NormalTok{(}\AttributeTok{component.autoscale =} \ConstantTok{FALSE}\NormalTok{)}
\NormalTok{NSmap}
\end{Highlighting}
\end{Shaded}

\pandocbounded{\includegraphics[keepaspectratio]{Day23_Lesson_files/figure-pdf/map simple-1.pdf}}

After a reasonably simple map, we need to add a few more parameters and
arguments to further customise the figure further. The command
\texttt{tm\_credits} allows us to add in text relating to data sources
and other items using the \texttt{\textbackslash{}n} adds a new line and
the position it where it needs to go, so the center and bottom of the
figure. I have also chosen a specific text \texttt{fontface} of bold and
\texttt{fontfamily} of mono and added a background that is black using
\texttt{bg.color} with white text using \texttt{color="white"} and
again, the \texttt{size} parameter changes the size of the text. Our
final tweak is to use the \texttt{tm\_layout} command to adjust the
background colour of the map and drop the frame, using \texttt{frame=F}
to get rid of the standard box around the plot.

\begin{Shaded}
\begin{Highlighting}[]
\NormalTok{NSmap }\OtherTok{\textless{}{-}}
  \FunctionTok{tm\_shape}\NormalTok{(UK) }\SpecialCharTok{+}
  \FunctionTok{tm\_polygons}\NormalTok{(}\StringTok{"DIVIDE"}\NormalTok{,}
              \AttributeTok{fill.scale =} \FunctionTok{tm\_scale}\NormalTok{(}\AttributeTok{values =}\NormalTok{ NS\_palette),}
              \AttributeTok{col=}\ConstantTok{NULL}\NormalTok{, }\AttributeTok{fill.legend =} \FunctionTok{tm\_legend\_hide}\NormalTok{())}\SpecialCharTok{+}
  \FunctionTok{tm\_title}\NormalTok{(}\AttributeTok{text=}\StringTok{"North{-}South Divide"}\NormalTok{, }\AttributeTok{size =} \FloatTok{1.5}\NormalTok{,}
           \AttributeTok{frame=}\NormalTok{F, }\AttributeTok{fontface =} \StringTok{"bold"}\NormalTok{, }\AttributeTok{fontfamily =} \StringTok{"mono"}\NormalTok{,}
           \AttributeTok{color =} \StringTok{"white"}\NormalTok{, }\AttributeTok{bg =} \ConstantTok{TRUE}\NormalTok{, }\AttributeTok{bg.color =} \StringTok{"black"}\NormalTok{, }
           \AttributeTok{position =} \FunctionTok{tm\_pos\_out}\NormalTok{(}\StringTok{"center"}\NormalTok{, }\StringTok{"top"}\NormalTok{, }\AttributeTok{pos.h =} \StringTok{"center"}\NormalTok{) )}\SpecialCharTok{+}
  \FunctionTok{tm\_layout}\NormalTok{(}\AttributeTok{bg.color =} \StringTok{"aliceblue"}\NormalTok{, }\AttributeTok{Frame=}\NormalTok{F)  }\SpecialCharTok{+}
  \FunctionTok{tm\_credits}\NormalTok{(}\AttributeTok{text =}\StringTok{"Created by Prof. Malcolm Campbell }\SpecialCharTok{\textbackslash{}n}\StringTok{Data: rnaturalearth boundaries"}\NormalTok{,}
             \AttributeTok{position =} \FunctionTok{tm\_pos\_out}\NormalTok{(}\StringTok{"center"}\NormalTok{, }\StringTok{"bottom"}\NormalTok{, }\AttributeTok{pos.h =} \StringTok{"center"}\NormalTok{),}
             \AttributeTok{size=}\FloatTok{0.75}\NormalTok{, }\AttributeTok{fontface =} \StringTok{"bold"}\NormalTok{, }\AttributeTok{fontfamily =} \StringTok{"mono"}\NormalTok{,}
             \AttributeTok{color =} \StringTok{"white"}\NormalTok{, }\AttributeTok{bg =} \ConstantTok{TRUE}\NormalTok{, }\AttributeTok{bg.color =} \StringTok{"black"}\NormalTok{, ) }\SpecialCharTok{+}
  \FunctionTok{tmap\_options}\NormalTok{(}\AttributeTok{component.autoscale =} \ConstantTok{FALSE}\NormalTok{)}
\NormalTok{NSmap}
\end{Highlighting}
\end{Shaded}

\pandocbounded{\includegraphics[keepaspectratio]{Day23_Lesson_files/figure-pdf/map it-1.pdf}}

Lastly, we might want to save our map and share it using the
\texttt{tmap\_save} command. We need to put in our map object, called
NSmap, choose a filename and format, for example .png alongside a
height, width and dots per inch (dpi) value.

\begin{Shaded}
\begin{Highlighting}[]
\FunctionTok{tmap\_save}\NormalTok{(NSmap, }\AttributeTok{filename =} \StringTok{"Day22\_NS.png"}\NormalTok{, }\AttributeTok{height =} \DecValTok{12}\NormalTok{, }\AttributeTok{width =} \DecValTok{12}\NormalTok{, }\AttributeTok{dpi=}\DecValTok{400}\NormalTok{)}
\end{Highlighting}
\end{Shaded}

There you have it. A `simple' map.




\end{document}
